\section*{Bab \ref{ch:g10:cells-common-unit} --- Sel: Satuan Bersama Kehidupan}

\begin{solution}{exo:g10:cells-common-unit:1}
Setiap makhluk hidup terbuat dari satu \omterm{def:g10:cells-common-unit:cell}{sel} atau lebih; \omterm{def:g10:cells-common-unit:cell}{sel} adalah
satuan dasar struktur dan fungsi; setiap \omterm{def:g10:cells-common-unit:cell}{sel} berasal dari \omterm{def:g10:cells-common-unit:cell}{sel} yang
sudah ada sebelumnya lewat pembelahan.
\end{solution}

\begin{solution}{exo:g10:cells-common-unit:2}
\qty{50}{\micro m}; \qty{2.5}{\micro m}; \qty{0.007}{mm}.
\end{solution}

\begin{solution}{exo:g10:cells-common-unit:3}
$36/600 = \qty{0.06}{mm} = \qty{60}{\micro m}$.
\end{solution}

\begin{solution}{exo:g10:cells-common-unit:4}
Hanya tumbuhan: \omterm{def:g10:cells-common-unit:organelle}{dinding sel}, \omterm{def:g10:cells-common-unit:organelle}{kloroplas}, \omterm{def:g10:cells-common-unit:organelle}{vakuola} pusat yang besar.
Setiap \omterm{def:g10:cells-common-unit:cell}{sel}: \omterm{def:g10:cells-common-unit:cell}{membran plasma} (juga \omterm{def:g10:cells-common-unit:cell}{sitoplasma} dan \omterm{def:g10:cells-common-unit:organelle}{ribosom}).
\end{solution}

\begin{solution}{exo:g10:cells-common-unit:5}
\omterm{def:g10:cells-common-unit:prokaryote}{Sel eukariotik} menyimpan DNA-nya di dalam \omterm{def:g10:cells-common-unit:organelle}{inti sel} dan punya \omterm{def:g10:cells-common-unit:organelle}{organel}
bermembran (\omterm{def:g10:cells-common-unit:cell}{sel} pipi, ragi); \omterm{def:g10:cells-common-unit:prokaryote}{sel prokariotik} tidak punya keduanya,
DNA-nya terletak bebas dalam \omterm{def:g10:cells-common-unit:cell}{sitoplasma}, dan ukurannya sekitar
\qty{1}{\micro m} (bakteri seperti \emph{E. coli}).
\end{solution}

\begin{solution}{exo:g10:cells-common-unit:6}
Panjang sebenarnya $6 \times 5/15 = \qty{2}{\micro m}$. \omterm{def:g10:cells-common-unit:magnification}{Perbesarannya}:
batangnya \qty{15}{mm} untuk \qty{5}{\micro m} $= \qty{0.005}{mm}$, jadi
$15/0.005 = 3000$: $\times 3000$.
\end{solution}

\begin{solution}{exo:g10:cells-common-unit:7}
Satu \omterm{def:g10:cells-common-unit:organelle}{ribosom} (\qty{25}{nm}) sepuluh kali lebih kecil daripada \omterm{def:g10:cells-common-unit:magnification}{daya
pisah} mikroskop cahaya (\qty{0.2}{\micro m} $= \qty{200}{nm}$): dua
\omterm{def:g10:cells-common-unit:organelle}{ribosom} yang bertetangga, atau satu \omterm{def:g10:cells-common-unit:organelle}{ribosom} dan sekelilingnya, tidak
dapat dibedakan. \omterm{def:g10:cells-common-unit:magnification}{Perbesaran} yang lebih besar hanya memperbesar kaburnya.
\end{solution}

\begin{solution}{exo:g10:cells-common-unit:8}
Butiran itu adalah pati, yang disimpan di dalam \omterm{def:g10:cells-common-unit:cell}{selnya} sebagai butiran
padat (pati tidak larut). Cadangan tumbuhan disimpan di dalam \omterm{def:g10:cells-common-unit:cell}{selnya},
bukan di suatu ruang di antara \omterm{def:g10:cells-common-unit:cell}{selnya}.
\end{solution}

\begin{solution}{exo:g10:cells-common-unit:9}
Sisik bawang bombai tumbuh di bawah tanah, dalam gelap, dan menyimpan
cadangan: \omterm{def:g10:cells-common-unit:organelle}{kloroplas} tidak akan berguna di sana. Daun lumut hidup dalam
cahaya dan membuat makanannya lewat fotosintesis, sehingga \omterm{def:g10:cells-common-unit:cell}{selnya} padat
dengan \omterm{def:g10:cells-common-unit:organelle}{kloroplas}. Rancangan sama, perlengkapan berbeda.
\end{solution}

\begin{solution}{exo:g10:cells-common-unit:10}
Sisi \qty{1}{\micro m}: permukaan \qty{6}{\micro m^2}, volume
\qty{1}{\micro m^3}, nisbahnya 6 tiap mikrometer. Sisi \qty{20}{\micro m}:
permukaan \qty{2400}{\micro m^2}, volume \qty{8000}{\micro m^3},
nisbahnya 0.3. Bakteri punya membran dua puluh kali lebih banyak tiap
satuan isi dan hidup dengan peresapan saja; \omterm{def:g10:cells-common-unit:cell}{sel} yang lebih besar
memerlukan pengangkutan dalam dan \omterm{def:g10:cells-common-unit:organelle}{organel}.
\end{solution}

\begin{solution}{exo:g10:cells-common-unit:11}
\qty{2}{h} berarti 6 pembelahan: $2^6 = 64$ \omterm{def:g10:cells-common-unit:cell}{sel}. \qty{4}{h} berarti 12:
$2^{12} = 4096$.
\end{solution}

\begin{solution}{exo:g10:cells-common-unit:12}
Ia punya membran, \omterm{def:g10:cells-common-unit:cell}{sitoplasma} serta menerima dan mengeluarkan zat,
sehingga ia memenuhi sebagian besar definisinya, tetapi ia tidak dapat
membelah dan tidak dapat memperbarui \omterm{def:g10:chemistry-of-life:families}{proteinnya}: ia adalah \omterm{def:g10:cells-common-unit:cell}{sel} yang
telah melepaskan sebagian programnya demi ruang bagi hemoglobin, dan
ia dibuat serta diganti oleh \omterm{def:g10:cells-common-unit:cell}{sel} yang utuh. Definisinya melukiskan
keadaan umum; \omterm{def:g10:cells-common-unit:cell}{sel} darah merah adalah bentuk khususnya yang terakhir.
\end{solution}

\begin{solution}{exo:g10:cells-common-unit:13}
Gabus adalah jaringan mati: Hooke melihat dinding selulosa \omterm{def:g10:cells-common-unit:cell}{sel} yang isi
hidupnya sudah lenyap --- kotak kosong. \omterm{def:g10:cells-common-unit:organelle}{Inti sel} hanya ada dalam
\omterm{def:g10:cells-common-unit:cell}{sitoplasma} \omterm{def:g10:cells-common-unit:cell}{sel} yang hidup, dan lagi pula \omterm{def:g10:cells-common-unit:magnification}{daya pisah} mikroskopnya serta
tidak adanya zat pewarna akan menyembunyikannya.
\end{solution}

\begin{solution}{exo:g10:cells-common-unit:14}
Virus tidak punya \omterm{def:g10:cells-common-unit:cell}{sitoplasma}, tidak punya membran buatannya sendiri,
tidak punya metabolisme, dan tidak dapat membelah: ia gagal memenuhi
definisi \omterm{def:g10:cells-common-unit:cell}{sel}. Menurut teori \omterm{def:g10:cells-common-unit:cell}{sel} ia tidak hidup sendirian; ia adalah
partikel yang meminjam permesinan \omterm{def:g10:cells-common-unit:cell}{sel} untuk membuat salinan dirinya.
Menyebutnya ``hidup'' atau tidak adalah soal definisi; bahwa ia bukan
\omterm{def:g10:cells-common-unit:cell}{sel} bukanlah soal definisi.
\end{solution}

\begin{solution}{exo:g10:cells-common-unit:15}
Volumenya $\pi \times (25\,\unit{\micro m})^2 \times 30000\,\unit{\micro m}
\approx \num{5.9e7}\,\unit{\micro m^3}$, sekitar 7000 kali
\qty{8000}{\micro m^3} kubusnya. Karena berupa tabung yang tipis, tidak
ada titik serat itu yang lebih dari \qty{25}{\micro m} dari membrannya,
sehingga pertukarannya tetap cepat; dan tiap \omterm{def:g10:cells-common-unit:organelle}{inti sel} hanya memerintah
bentangnya sendiri pada serat itu, sehingga tidak ada \omterm{def:g10:cells-common-unit:organelle}{inti sel} yang
harus melayani volume yang lebih besar daripada volume \omterm{def:g10:cells-common-unit:cell}{sel} biasa.
\end{solution}

\begin{solution}{pb:g10:cells-common-unit:1}
\textbf{1.} $1.8/9 = \qty{0.2}{mm} = \qty{200}{\micro m}$.

\textbf{2.} $1.8/4 = \qty{0.45}{mm}$; sekitar dua \omterm{def:g10:cells-common-unit:cell}{sel} muat melintasinya.

\textbf{3.} $28/0.2 = 140$: gambarnya $\times 140$.

\textbf{4.} $4/140 \approx \qty{0.029}{mm}$, sekitar \qty{30}{\micro m}.

\textbf{5.} \omterm{def:g10:cells-common-unit:organelle}{Dinding sel}, \omterm{def:g10:cells-common-unit:cell}{membran plasma} (yang menempel pada dinding),
\omterm{def:g10:cells-common-unit:cell}{sitoplasma}, \omterm{def:g10:cells-common-unit:organelle}{inti sel}, \omterm{def:g10:cells-common-unit:organelle}{vakuola}. Hadir tetapi tak terlihat pada
$\times 400$: \omterm{def:g10:cells-common-unit:organelle}{ribosomnya} (dan \omterm{def:g10:cells-common-unit:organelle}{mitokondria} ada di batasnya).

\textbf{6.} \omterm{def:g10:cells-common-unit:cell}{Sel} itu berasal dari permukaan lapisan yang tebalnya
beberapa \omterm{def:g10:cells-common-unit:cell}{sel} dan terus-menerus tergosok: \omterm{def:g10:cells-common-unit:cell}{sel} terluarnya memipih dan
terlepas, dan itulah sebabnya kapas bertangkai memungutnya.

\textbf{7.} \qty{0.24}{mm} hanya sedikit di atas batas
\qty{0.1}{mm} mata telanjang: sebuah bintik yang nyaris terlihat. Pada
$\times 100$ ia melintasi $0.24/1.8 \approx 13\%$ bidang pandangnya,
pada $\times 400$ lebih dari separuh; tidak ada lensa objektif yang
tersedia memberi tepat seperempat.

\textbf{8.} Bakteri. Mikroskop elektron akan memperlihatkan dinding,
membran, \omterm{def:g10:cells-common-unit:organelle}{ribosom} dan daerah DNA-nya, serta tidak adanya \omterm{def:g10:cells-common-unit:organelle}{inti sel}.

\textbf{9.} \omterm{def:g10:cells-common-unit:cell}{Sel} pipi mengerjakan satu tugas (melindungi sebuah
permukaan) dalam tubuh tempat \omterm{def:g10:cells-common-unit:cell}{sel} lain memberinya makan, memasok
oksigen dan menyelaraskannya; paramesium tidak punya siapa pun dan
harus makan, bergerak, mengindra dan membelah sebagai \omterm{def:g10:cells-common-unit:cell}{sel} tunggal.

\textbf{10.} Bagian luarnya kini lebih asin daripada \omterm{def:g10:cells-common-unit:cell}{sitoplasmanya},
sehingga air meninggalkan \omterm{def:g10:cells-common-unit:cell}{selnya} melintasi membran; \omterm{def:g10:cells-common-unit:cell}{selnya} kehilangan
volume dan mengerut.

\textbf{11.} $20^3 = \qty{8000}{\micro m^3} = \num{8e-6}\,\unit{mm^3}$.

\textbf{12.} $6 \times 20^2 = \qty{2400}{\micro m^2}$; nisbahnya
$2400/8000 = 0.3$ tiap mikrometer.

\textbf{13.} Permukaan \qty{6}{\micro m^2}, volume \qty{1}{\micro m^3},
nisbahnya 6: dua puluh kali lebih besar.

\textbf{14.} $\qty{8000}{\micro m^3} = \num{8e-9}\,\unit{cm^3}$, jadi
sekitar $\num{8e-9}\,\unit{g} = \qty{8}{ng}$.

\textbf{15.} $\pi \times 3.5^2 \times 2 \approx \qty{77}{\micro m^3}$,
sekitar seperseratus kubusnya.

\textbf{16.} $\qty{60}{L} = \num{6e16}\,\unit{\micro m^3}$; dibagi
8000: sekitar $\num{7.5e12}$ \omterm{def:g10:cells-common-unit:cell}{sel}.

\textbf{17.} Dua pertiganya: sekitar $\num{5e12}$.

\textbf{18.} $\num{2.5e13} \times 77 \approx \num{1.9e15}\,\unit{\micro
m^3} \approx \qty{1.9}{L}$. \omterm{def:g10:cells-common-unit:cell}{Sel} darahnya: $0.45 \times 4.5 \approx
\qty{2.0}{L}$ --- sesuai.

\textbf{19.} Sekitar $\num{5e12}$ \omterm{def:g10:cells-common-unit:cell}{sel} besar ditambah $\num{2.5e13}$ \omterm{def:g10:cells-common-unit:cell}{sel}
merah: sekitar $\num{3e13}$, beberapa puluh triliun. Ukuran \omterm{def:g10:cells-common-unit:cell}{sel}
merentang sampai faktor seratus dan campuran jenisnya berubah-ubah,
sehingga penghitungan apa pun adalah taksiran sampai faktor dua atau tiga.

\textbf{20.} Kira-kira tiga puluh triliun ($\num{3e13}$) \omterm{def:g10:cells-common-unit:cell}{sel}, semuanya
dihasilkan lewat pembelahan \omterm{def:g10:cells-common-unit:cell}{sel} berturut-turut dari satu \omterm{def:g10:cells-common-unit:cell}{sel} awal,
yaitu \omterm{def:g10:cells-common-unit:cell}{sel} telur yang telah dibuahi.
\end{solution}
